\subsection{Learning theoretic analysis}\label{subsec:learning-theoretic-analysis}
\subsubsection{Effective task learning}\label{subsec:effective-learning}
In this section, we show that we can effectively learn the task functional under mild conditions on the quantum resources used. Let our dataset of windows be $\gD = \{(\vx_1, y_1), \dots, (\vx_N, y_N) \}$. We argue as follows that, upon using a quantum reservoir computer with a number of qubits as low as $n = \Omega(\log(w N))$, we can, with high probability, achieve effective empirical learning. In this setting, we show in the following theorem that, in an idealized noiseless setting (as opposed to the noisy feature maps of Section~\ref{subsubsec:efficient-msmt-scheme}), we obtain effective learning of the task at hand. The proof of such result can be found in Appendix~\ref{subapp:proof-thm-effective-learning}.

\begin{theorem}[Effective learning]\label{thm:effective-learning}
    Suppose a design of a quantum reservoir as described in Section~\ref{subsec:QR-embedding} with a number of $R \geq 1$ sub-reservoirs, and let our readout scheme follow the kernel-based one described in Section~\ref{subsec:K-read}. Let $\varepsilon_{\pr}, \delta_{\pr}  \in (0,1)$ be, respectively, an error tolerance and a failure probability parameter. Then, with a number of qubits as low as $n = \Omega\pth{\varepsilon_{\pr}^{-2} \log\pth{\tfrac{w N}{\delta_{\pr}}}}$, with probability at least $1-\delta_{\pr}$, there exists a value $\Lambda^\star > 0$ such that for any constraint $\Lambda \leq \Lambda^\star$, we have
    \begin{equation}\label{eq:empirical-risk-and-RKHS-norm}
        \min_{h\in \gHkp(\Lambda)} \hat{R}_N(H^T_h) = O\pth{\big(1 - \tfrac{\Lambda}{\Lambda^\star}\big)^2}.
    \end{equation}
\end{theorem}

Theorem~\ref{thm:effective-learning} can be read as a capacity statement for the end-to-end functional $H^T_h$: for a fixed quantum featurizer, enlarging the RKHS norm constraint decreases the achievable training loss. This is made explicit in~\cref{eq:empirical-risk-and-RKHS-norm} by exhibiting a budget threshold (which depends on the Gram matrix) such that, once $\Lambda$ is relaxed enough to approach $\Lambda^\star$, the empirical risk can be driven arbitrarily close to zero. Therefore, under the stated qubit scaling (which leads to a JL projection preserving the finite window set's geometry with high probability), the kernel readout becomes expressive enough to interpolate the dataset under sufficiently weak regularization. The interpretation of this result should be solely that there exists an effective interpolation regime.


\subsubsection{Generalization analysis on weakly-dependent data}\label{subsubsec:generalization}
Though Theorem~\ref{thm:effective-learning} shows that empirical fit is achievable with modest quantum resources, it says nothing about how the empirical risk deviates from the true risk $R(H^T_h)$ as defined in~\cref{eq:risk-def}, i.e., generalization is not guaranteed. Characterizing the generalization gap $\abs{R(H^T_h) - \hat{R}_N(H^T_h)}$ is, in general, a non-trivial endeavour for time series data, as we do not necessarily have the required i.i.d.\ property for the window samples of our hookups dataset $\gD$ (between-window dependence may occur), which is often assumed by standard learning-theoretic analyses.

In our treatment, we provide a generalization characterization for the case where our data are weakly dependent and are sampled from a $\beta$-mixing process $\rmIO = ((X_t, Y_t) : t \in \sZ_-)$. In essence, a $\beta$-mixing process means that the ``far past'' of the sequence and the ``far future'' become more and more independent as we increase the temporal gap (formal definitions can be found in Appendix~\ref{subsec:weakly-dependent-data}). This property can be leveraged to enable coupling/blocking arguments that lead to generalization bounds similar to the i.i.d.\ case, with an explicit dependence-related penalty~\cite{mohri2008rademacher}. It is worth mentioning that this weak-dependence assumption encompasses a broad class of time series models used in practice, which are $\beta$-mixing under verifiable stability/ergodicity conditions, making our analysis applicable well beyond the i.i.d.\ setting~\cite{dedecker2007,Doukhan_1994,Carrasco2002}.

We now describe how we construct a dataset that leads to our generalization guarantee. We observe a stationary input/output process $\rmIO$ (see Section~\ref{subsec:learning-temporal-data} for details). Then, a window $\rvx_i$ is constructed as the realization of the sub-process $(X_{t_i - w + 1}, \ldots, X_{t_i})$ and is assigned a label $y_i$, the realization of $Y_{t_i}$. When windows overlap, dependence arises; hence, we instead select indices with stride $s = w + g$, where $g$ is a controllable gap. Increasing $g$ increases, in turn, the independence between windows for our $\beta$-mixing process. Sampling $N$ windows while respecting stride $s$ between consecutive windows is how we construct our dataset $\gD = \{(\rvx_1, y_1), \ldots, (\rvx_N, y_N)\}$, where we assume $t_1 > t_2 > \cdots > t_N$ with $t_i - t_{i+1} = s$. We now present our generalization result for the $\beta$-mixing case the proof of which can be found in Appendix~\ref{subsubsec:proof-thm-generalization}.


\begin{theorem}[Generalization on weakly-dependent data]\label{thm:generalization}
Consider $\rmIO = ((X_t,Y_t):t)$ as a stationary $\beta$-mixing process with bounded outputs
$|Y_t|\leq \Upsilon_Y$. Fix a window length $w$ and a gap $g$, and construct a strided windows
dataset $\gD=\{(\rvx_i,y_i)\}_{i=1}^N$ with stride $s=w+g$ (as in the paragraph above). Assume $N$ is even and write $N=2\mu$. Let $T$ be a QR featurizer composed of $R$ sub-reservoirs with contraction values $\{\lambda_r\}_{r=1}^R$, and let $\kappa$ be the unit-variance Mat\'ern kernel used in Section~\ref{subsec:K-read}
(with parameters $\nu>1$ and $\xi>0$). Let $\gO$ be the set of measured $k$-local observables.
Define $\lambda_\star := \max_{r\in[R]}\lambda_r$.

Fix any $\delta\in(0,1)$ such that
\begin{equation}
    \delta > 4(\mu-1)\beta_{\rmIO}(g),
    \qquad
    \delta' := \delta - 4(\mu-1)\beta_{\rmIO}(g) \;>\; 0.
\end{equation}
Then, with probability at least $1-\delta$, the following holds \emph{simultaneously for all}
$h\in \gHkp(\Lambda)$:
\begin{equation}\label{eq:bound-generalization}
\begin{aligned}
\abs{R(H^T_h) - \hat{R}_N(H^T_h)}
\;\leq\;
&\underbrace{4\sqrt{2}\,\frac{\Lambda(\Lambda+\Upsilon_Y)}{\sqrt{N}}}_{\text{(Rademacher term)}}
\;+\;
\underbrace{3(\Lambda+\Upsilon_Y)^2\frac{\sqrt{\log(4/\delta')}}{\sqrt{N}}}_{\text{(mixing penalty)}}\\
&\;+\;
\underbrace{\frac{4\Lambda(\Lambda+\Upsilon_Y)}{\xi}\sqrt{\frac{\nu R|\gO|}{\nu-1}}\,
\lambda_\star^{\,w}}_{\text{(window truncation via contraction)}}.
\end{aligned}
\end{equation}
\end{theorem}

Theorem~\ref{thm:generalization} provides a uniform high-probability bound on the generalization gap $|R(H^T_h) - \hat{R}_N(H^T_h)|$ for all readouts $h$ within the RKHS ball $\gHkp(\Lambda)$ trained on a strided windows dataset. The bound decomposes into three terms. First, a standard Rademacher complexity term of order $1/\sqrt{N}$, which penalizes model richness. Second, a dependence penalty that is also $O(1/\sqrt{N})$ but whose confidence level is effectively reduced by temporal dependence through $\delta' = \delta - 4(\mu - 1)\beta_{\rmIO}(g)$; hence the need to choose a gap $g$ large enough to reduce $\beta_{\rmIO}(g)$. Third, a fading-memory remainder that decays geometrically with $\lambda_\star^w$, which penalizes the fact that the reservoir forgets information beyond the window horizon. Consequently, for a fixed design, increasing the number of windows $N$ drives the first two terms to zero, while choosing $w$ moderately large (or ensuring stronger contraction) makes the truncation term negligible. In practice, the gap $g$ controls statistical dependence between windows, while $w$ controls the approximation to the reservoir's infinite-memory dynamics, thereby yielding a clear bias--dependence--variance tradeoff.



% The right-hand side of~\eqref{eq:bound-generalizatio} consists of three parts: (i) an estimation term $O(1/\sqrt{N})$ captured by the empirical Rademacher complexity, (ii) a concentration term that depends on the confidence level and effective sample size, and (iii) a geometric decay term $G(\kappa,\Lambda,T)$ that reflects the fading-memory behavior induced by the contraction parameters and window length. These terms show that, for fixed design and resource choices, the generalization error can be made small by increasing the number of windows $N$ while controlling dependence through the stride $s=w+g$.
% The utility of this theorem is that, for any fixed configuration $(\ell,\kappa,w,g,\Lambda,T)$, the generalization error bound can be made small by increasing the number of windows $N$. This behavior is empirically illustrated in Section~4.2.2, where the test error decreases as $N$ grows.

%The utility of this theorem arises from the fact that provided a fixed configuration $(\ell, \kappa, w, g, \Lambda, T)$, we can push the generalization error bound towards a small value by working with a large sample set $N$. This is what we validate empirically in our experiment of Section~\ref{subsubsec:generalization-validation}. \TBDcomment{[See if more discussion needed here!]}
